% ===================================================================
%  CUADRO N.º 13 — El Hotel. — El Restaurante. — El Café.
%  Traduction 2026 d'apres texto/io, controlee sur texto/fr et
%  texto/en. Consigne : voir texto/es/10-cuadro-01.tex.
%
%  Deux notes, marquees toutes deux « (*) », et deux appels : celui du
%  bloc 01-2 (la chambre) et celui du bloc 09-1 (le menu). L'ordre les
%  departage.
% ===================================================================

%%K t13-noto-1 noto
\VUnotes{143.2mm}{%
\textit{\VUgras{(*) Para los detalles del dormitorio, véase el cuadro n.º 6.}}%
}

%%K t13-apar-1 apar
\VUsaut{3.0mm}
\VUtitre{15.0pt}{60}{CUADRO N.º 13}
\VUsaut{1.6mm}
\VUpk{10.2pt}{40}{El Hotel. --- El Restaurante.}
\VUsaut{2.0mm}
\VUpk{10.2pt}{40}{El Café.}
\VUsaut{3.4mm}

%%K t13-01-1 p
1. --- Habiendo llegado a Royan ayer por la tarde, me alojé en el
\textit{Hotel del Océano}, que me había recomendado un amigo.

%%K t13-01-2 p
Me dieron una hermosa \VUgras{habitación} \textsuperscript{(2)} en el segundo
\VUgras{piso} \textsuperscript{(1)}, donde dormí muy bien (*).

%%K t13-02-1 p
2. --- Esta mañana, al salir de mi habitación, donde la \VUgras{camarera}
\textsuperscript{(3)} me había traído el desayuno, entré en el \VUgras{salón de
fumar} \textsuperscript{(5)}, que sirve también de \VUgras{sala de lectura}
\textsuperscript{(4)}, para hojear los \VUgras{periódicos} \textsuperscript{(6)} mientras
fumaba un \VUgras{cigarrillo} \textsuperscript{(8)}. Cuando terminé de leer, bajé en
el \VUgras{ascensor} \textsuperscript{(9)} y salí a dar un paseo por la ciudad.

%%K t13-03-1 p
3. --- Como se me había olvidado preguntar al \VUgras{recepcionista}
\textsuperscript{(12)} del hotel a qué hora se servían las comidas en la
\VUgras{mesa redonda} \textsuperscript{(11)}, volví temprano y aproveché para darme
un baño en el \VUgras{cuarto de baño} \textsuperscript{(10)} del hotel. Después fui
a la \VUgras{caja} \textsuperscript{(15)} para telefonear a uno de mis amigos. Pero
el \VUgras{teléfono} \textsuperscript{(13)} estaba ocupado; al ver mi apuro, la
\VUgras{cajera} \textsuperscript{(14)}, que estaba anotando una \VUgras{cuenta}
\textsuperscript{(17)} en el \VUgras{registro de viajeros} \textsuperscript{(16)}, pidió al
\VUgras{conserje} \textsuperscript{(18)} que mandara al \VUgras{botones}
\textsuperscript{(19)} a hacer mi recado en su \VUgras{bicicleta} \textsuperscript{(20)}.

%%K t13-04-1 p
4. --- Le di las gracias y salí a dar una propina al \VUgras{mozo de cuerda}
\textsuperscript{(24)} (el hombre para todo), que había traído de la estación en
su \VUgras{carretilla} \textsuperscript{(25)} mi \VUgras{sombrerera}
\textsuperscript{(26)}, mi \VUgras{maleta} \textsuperscript{(27)} y mi \VUgras{bolsa de
viaje} \textsuperscript{(28)}. El \VUgras{cartero} \textsuperscript{(21)}, que llegaba en
ese momento, me entregó una \VUgras{carta} \textsuperscript{(22)}.

%%K t13-05-1 p
5. --- Me quedé un rato más en el umbral, junto al \VUgras{buzón}
\textsuperscript{(23)}, mirando el movimiento de la calle. El \VUgras{conductor}
\textsuperscript{(30)} del \VUgras{ómnibus} \textsuperscript{(29)} cargaba en su vehículo
el \VUgras{baúl} \textsuperscript{(31)} de un \VUgras{viajero} \textsuperscript{(32)} a quien
el \VUgras{hotelero} \textsuperscript{(33)} decía adiós. Un joven
\VUgras{telegrafista} \textsuperscript{(35)} traía, sin ninguna prisa, un
\VUgras{telegrama} \textsuperscript{(34)} para uno de los huéspedes del hotel; un
\VUgras{turista} \textsuperscript{(36)}, con su \VUgras{cámara fotográfica}
\textsuperscript{(38)} colgada del hombro, consultaba su \VUgras{guía}
\textsuperscript{(37)}.

%%K t13-tit-1 sub
\VUsaut{4.0mm}
\VUpk{10.2pt}{40}{La hora del almuerzo.}
\VUsaut{3.0mm}

%%K t13-06-1 p
6. --- Delante de la verja, un plácido \VUgras{recadero} \textsuperscript{(39)}
dormitaba entre sus \VUgras{ganchos de carga} \textsuperscript{(40)} y su
\VUgras{caja de limpiabotas} \textsuperscript{(41)}, mientras una joven
\VUgras{lavandera} \textsuperscript{(42)} que llevaba un \VUgras{cesto}
\textsuperscript{(43)} de ropa se reía del aire solemne del \VUgras{maître}
\textsuperscript{(44)} plantado junto a la puerta.

%%K t13-07-1 p
7. --- La vista del \VUgras{cocinero} \textsuperscript{(45)}, que había salido de su
\VUgras{cocina} \textsuperscript{(47)} del \VUgras{sótano} \textsuperscript{(48)} para mandar
a un \VUgras{pinche} \textsuperscript{(46)} a un recado, me recordó que se acercaba
la hora del almuerzo, y fui al restaurante al otro lado del patio, donde
un \VUgras{automovilista} \textsuperscript{(50)} aparcaba su \VUgras{automóvil}
\textsuperscript{(49)} mientras un \VUgras{mecánico} \textsuperscript{(52)} reparaba,
siguiendo las indicaciones del \VUgras{ciclista} \textsuperscript{(53)}, un
\VUgras{triciclo de gasolina} \textsuperscript{(51)} que acababa de tener un
accidente.

%%K t13-08-1 p
8. --- Pregunté a un joven \VUgras{groom} \textsuperscript{(54)} que llevaba
\VUgras{jarras de cerveza} \textsuperscript{(55)} si la mesa redonda estaba bajo la
veranda, y como me dijo que sí, me senté a una de las mesas y me sirvieron
una comida excelente.

%%K t13-noto-2 noto
\VUnotes{143.2mm}{%
\textit{\VUgras{(*) Con ayuda de la lista de platos que da el vocabulario,
el maestro puede variar fácilmente el menú siguiente.}}%
}

%%K t13-tit-2 sub
\VUsaut{4.0mm}
\VUpk{10.2pt}{40}{Un almuerzo excelente.}
\VUsaut{3.0mm}

%%K t13-09-1 p
9. --- El camarero me trajo el menú del almuerzo (*) y la carta de vinos.
Elegí y pedí \VUgras{gambas} con \VUgras{mantequilla} de
\VUgras{entremés}, y, de primer plato, \VUgras{callos a la moda de Caen}. No
tomé \VUgras{pescado}, pero pedí una ración de \VUgras{pierna} de cordero y,
de \VUgras{verdura}, \VUgras{espinacas} a la crema. Después, como ninguno de
los \VUgras{platos dulces} me apetecía, hice traer un surtido de postres:
\VUgras{queso}, \VUgras{fruta} y \VUgras{galletas}. Añadí un excelente
\VUgras{vino} de Saint-Émilion y, muy satisfecho de mi comida, me levanté
de la mesa después de dar al \VUgras{camarero} \textsuperscript{(57)} una buena
\VUgras{propina} \textsuperscript{(59)}.

%%K t13-10-1 p
10. --- Al ver que me disponía a marcharme, el \VUgras{gerente}
\textsuperscript{(60)} me propuso salir al balcón para admirar el espléndido
panorama del \VUgras{océano} \textsuperscript{(62)}. A lo lejos se distinguían
pequeñas \VUgras{barcas} \textsuperscript{(63)} de pesca, \VUgras{veleros}
\textsuperscript{(64)} y un enorme \VUgras{transatlántico} \textsuperscript{(65)} cuya
silueta negra se recortaba sobre las \VUgras{nubes} \textsuperscript{(67)} blancas
del \VUgras{horizonte} \textsuperscript{(66)}. Una brisa ligera agitaba la
\VUgras{bandera} \textsuperscript{(70)} plantada encima del \VUgras{rótulo}
\textsuperscript{(69)} del \VUgras{vestíbulo} \textsuperscript{(68)}.

%%K t13-tit-3 sub
\VUsaut{3.0mm}
\VUpk{10.2pt}{40}{Diversiones.}
\VUsaut{3.0mm}

%%K t13-11-1 p
11. --- El espectáculo era tan interesante que decidí subir al día
siguiente a la terraza del café, llevando conmigo unos \VUgras{gemelos}
\textsuperscript{(71)} o un \VUgras{catalejo} \textsuperscript{(72)}.

%%K t13-12-1 p
12. --- Al salir del balcón fui al café del hotel a tomar una
\VUgras{consumición} \textsuperscript{(73)} y a jugar una partida de
\VUgras{billar} \textsuperscript{(75)}. Atravesé la sala del café, donde muchos
clientes jugaban al \VUgras{ajedrez} \textsuperscript{(78)}, a las \VUgras{damas}
\textsuperscript{(79)}, al \VUgras{dominó} \textsuperscript{(80)}, al \VUgras{chaquete}
\textsuperscript{(81)} o a las \VUgras{cartas} \textsuperscript{(82)}, y subí derecho a la
\VUgras{sala de billar} \textsuperscript{(74)}.

%%K t13-13-1 p
13. --- Terminada la partida, bajé a la planta baja y pedí al camarero un
\VUgras{sobre} \textsuperscript{(84)}, \VUgras{papel} \textsuperscript{(83)}, un
\VUgras{sello} \textsuperscript{(85)}, una \VUgras{pluma} \textsuperscript{(87)} y
\VUgras{tinta} \textsuperscript{(86)} para escribir una carta.

%%K t13-13-2 p
Después llamé a un \VUgras{cochero} \textsuperscript{(92)} cuyo \VUgras{coche}
\textsuperscript{(88)} se había parado delante del café. Le pregunté su precio
por hora y por carrera y, una vez de acuerdo, me abrió la
\VUgras{portezuela} \textsuperscript{(89)} y me instaló dentro. Volvió a subir a su
\VUgras{pescante} \textsuperscript{(91)}, arreó vivamente a los caballos con su
\VUgras{látigo} \textsuperscript{(93)}, y allá fuimos a dar un paseo delicioso.
\VUsaut{6.0mm}
\VUfilet{22mm}
